\DocumentMetadata{
  lang=en-US,
  pdfversion=2.0,
  pdfstandard=ua-2,
  tagging=on,
  testphase={phase-III,math,title,table,tikz}
}
\documentclass[11pt]{article}
\usepackage[letterpaper,margin=0.55in]{geometry}
\usepackage{amsmath}
\usepackage{tikz}
\usetikzlibrary{arrows.meta,patterns,decorations.pathreplacing,calc}
\usepackage{xcolor}
\usepackage{multicol}
\usepackage{enumitem}
\usepackage{unicode-math}
\setmonofont{Latin Modern Mono}
\usepackage{hyperref}
\hypersetup{
  pdftitle={How the UV Sphere Is Built},
  pdfauthor={CS 4280},
  pdfsubject={CS 4280 sphere-shaded activity: UV sphere geometry},
  pdfkeywords={UV sphere, spherical coordinates, triangle mesh, normals, WebGPU},
  pdfdisplaydoctitle=true
}
\pagestyle{empty}
\setlength{\parindent}{0pt}
\setlength{\parskip}{3pt}
\setlength{\columnsep}{18pt}
\definecolor{thetac}{HTML}{C0392B}
\definecolor{phic}{HTML}{1F6FB2}
\definecolor{cellc}{HTML}{E08E0B}
\definecolor{keep}{HTML}{3A9D5D}
\newcommand{\code}[1]{\texttt{\small #1}}
\makeatletter
\renewcommand\section{\@startsection{section}{1}{0pt}{0pt}{-1em}{\LARGE\bfseries}}
\renewcommand\subsection{\@startsection{subsection}{2}{0pt}{4pt}{1pt}{\large\bfseries}}
\makeatother
\setcounter{secnumdepth}{0}
\newcommand{\step}[1]{\subsection{#1}}

\newcommand{\figonealt}{Three-dimensional diagram of a sphere of radius r centered at origin O, with the x, y, and z axes drawn; y points up. A point P on the sphere is joined to O by a radius. The angle theta is measured from the positive y axis (the north pole, theta = 0) down to the radius, and the angle phi is measured around the y axis starting from the positive z axis. Dashed helper lines show that P is at height r cos theta and at distance r sin theta from the y axis. Latitude and longitude grid lines cover the sphere, and one grid cell below the equator is shaded orange with corners a and b on its upper edge, c and d on its lower edge, and a diagonal from a to d. The south pole is at theta = pi.}
\begin{document}
\section{How the UV Sphere Is Built}\hfill{\small CS 4280 \,·\, \code{sphere-shaded} activity}\par
\vspace{2pt}\hrule\vspace{8pt}

\noindent
\begin{minipage}[c]{0.44\textwidth}\centering
\begin{tikzpicture}[alt={\figonealt},
  grid/.style={gray!45, line width=0.35pt},
  hgrid/.style={gray!45, dashed, line width=0.35pt},
  axis/.style={-{Stealth[length=5pt]}, line width=0.7pt},
  axisin/.style={gray!70, line width=0.4pt, dash pattern=on 1.5pt off 1.5pt},
  helper/.style={dashed, line width=0.6pt, black!70},
  radius/.style={line width=1pt, black},
  thetaarc/.style={thetac, line width=0.9pt}, thetacol/.style={thetac},
  phiarc/.style={phic, line width=0.9pt}, phicol/.style={phic},
  cellfill/.style={fill=cellc!30}, celldiag/.style={cellc, line width=0.5pt},
  cellpt/.style={cellc!80!black}, polecol/.style={black}]
\shade[ball color=gray!6, opacity=0.9] (0,0) circle (2.3);
\fill[cellfill] (-0.000,-0.711) -- (0.060,-0.710) -- (0.120,-0.710) -- (0.180,-0.709) -- (0.240,-0.707) -- (0.300,-0.705) -- (0.360,-0.702) -- (0.419,-0.699) -- (0.478,-0.695) -- (0.537,-0.691) -- (0.595,-0.686) -- (0.653,-0.681) -- (0.711,-0.676) -- (0.768,-0.670) -- (0.824,-0.663) -- (0.880,-0.657) -- (0.935,-0.649) -- (0.990,-0.641) -- (1.044,-0.633) -- (1.097,-0.625) -- (1.150,-0.615) -- (1.150,-0.615) -- (1.150,-0.673) -- (1.148,-0.729) -- (1.146,-0.785) -- (1.144,-0.841) -- (1.140,-0.896) -- (1.136,-0.950) -- (1.131,-1.004) -- (1.125,-1.057) -- (1.118,-1.109) -- (1.111,-1.161) -- (1.103,-1.211) -- (1.094,-1.261) -- (1.084,-1.310) -- (1.074,-1.359) -- (1.062,-1.406) -- (1.051,-1.452) -- (1.038,-1.497) -- (1.025,-1.542) -- (1.011,-1.585) -- (0.996,-1.627) -- (0.996,-1.627) -- (0.950,-1.635) -- (0.904,-1.642) -- (0.857,-1.649) -- (0.810,-1.656) -- (0.762,-1.662) -- (0.714,-1.668) -- (0.665,-1.674) -- (0.615,-1.679) -- (0.566,-1.684) -- (0.515,-1.688) -- (0.465,-1.692) -- (0.414,-1.696) -- (0.363,-1.699) -- (0.312,-1.702) -- (0.260,-1.704) -- (0.208,-1.706) -- (0.156,-1.707) -- (0.104,-1.708) -- (0.052,-1.709) -- (-0.000,-1.709) -- (-0.000,-1.709) -- (-0.000,-1.668) -- (-0.000,-1.626) -- (-0.000,-1.583) -- (-0.000,-1.539) -- (-0.000,-1.494) -- (-0.000,-1.447) -- (-0.000,-1.400) -- (-0.000,-1.352) -- (-0.000,-1.303) -- (-0.000,-1.253) -- (-0.000,-1.202) -- (-0.000,-1.150) -- (-0.000,-1.097) -- (-0.000,-1.044) -- (-0.000,-0.990) -- (-0.000,-0.935) -- (-0.000,-0.880) -- (-0.000,-0.824) -- (-0.000,-0.768) -- (-0.000,-0.711) -- cycle;
\draw[grid] (-0.575,1.587) -- (-0.540,1.581) -- (-0.504,1.575) -- (-0.468,1.570) -- (-0.431,1.565) -- (-0.393,1.561) -- (-0.355,1.557) -- (-0.317,1.553) -- (-0.278,1.550) -- (-0.239,1.547) -- (-0.200,1.545) -- (-0.160,1.543) -- (-0.120,1.541) -- (-0.080,1.540) -- (-0.040,1.539) -- (-0.000,1.539) -- (0.040,1.539) -- (0.080,1.540) -- (0.120,1.541) -- (0.160,1.543) -- (0.200,1.545) -- (0.239,1.547) -- (0.278,1.550) -- (0.317,1.553) -- (0.355,1.557) -- (0.393,1.561) -- (0.431,1.565) -- (0.468,1.570) -- (0.504,1.575) -- (0.540,1.581) -- (0.575,1.587) -- (0.609,1.593) -- (0.643,1.600) -- (0.676,1.607) -- (0.708,1.614) -- (0.739,1.622) -- (0.769,1.630) -- (0.799,1.639) -- (0.827,1.648) -- (0.855,1.657) -- (0.881,1.666) -- (0.906,1.676) -- (0.930,1.686) -- (0.953,1.696) -- (0.975,1.706) -- (0.996,1.717) -- (1.015,1.728) -- (1.034,1.739) -- (1.051,1.750) -- (1.066,1.761) -- (1.081,1.773) -- (1.094,1.785) -- (1.105,1.797) -- (1.116,1.808) -- (1.125,1.821) -- (1.133,1.833) -- (1.139,1.845) -- (1.144,1.857) -- (1.147,1.870) -- (1.149,1.882) -- (1.150,1.894) -- (1.149,1.907) -- (1.147,1.919) -- (1.144,1.932) -- (1.139,1.944) -- (1.133,1.956) -- (1.125,1.968) -- (1.116,1.980) -- (1.105,1.992) -- (1.094,2.004) -- (1.081,2.016) -- (1.066,2.028) -- (1.051,2.039) -- (1.034,2.050) -- (1.015,2.061) -- (0.996,2.072) -- (0.975,2.083) -- (0.953,2.093);
\draw[grid] (-0.930,2.103) -- (-0.953,2.093) -- (-0.975,2.083) -- (-0.996,2.072) -- (-1.015,2.061) -- (-1.034,2.050) -- (-1.051,2.039) -- (-1.066,2.028) -- (-1.081,2.016) -- (-1.094,2.004) -- (-1.105,1.992) -- (-1.116,1.980) -- (-1.125,1.968) -- (-1.133,1.956) -- (-1.139,1.944) -- (-1.144,1.932) -- (-1.147,1.919) -- (-1.149,1.907) -- (-1.150,1.894) -- (-1.149,1.882) -- (-1.147,1.870) -- (-1.144,1.857) -- (-1.139,1.845) -- (-1.133,1.833) -- (-1.125,1.821) -- (-1.116,1.808) -- (-1.105,1.797) -- (-1.094,1.785) -- (-1.081,1.773) -- (-1.066,1.761) -- (-1.051,1.750) -- (-1.034,1.739) -- (-1.015,1.728) -- (-0.996,1.717) -- (-0.975,1.706) -- (-0.953,1.696) -- (-0.930,1.686) -- (-0.906,1.676) -- (-0.881,1.666) -- (-0.855,1.657) -- (-0.827,1.648) -- (-0.799,1.639) -- (-0.769,1.630) -- (-0.739,1.622) -- (-0.708,1.614) -- (-0.676,1.607) -- (-0.643,1.600) -- (-0.609,1.593) -- (-0.575,1.587);
\draw[grid] (-0.996,0.561) -- (-0.935,0.550) -- (-0.873,0.541) -- (-0.810,0.531) -- (-0.746,0.523) -- (-0.681,0.515) -- (-0.616,0.508) -- (-0.549,0.502) -- (-0.482,0.497) -- (-0.414,0.492) -- (-0.346,0.488) -- (-0.277,0.484) -- (-0.208,0.482) -- (-0.139,0.480) -- (-0.070,0.479) -- (-0.000,0.478) -- (0.069,0.479) -- (0.139,0.480) -- (0.208,0.482) -- (0.277,0.484) -- (0.346,0.488) -- (0.414,0.492) -- (0.482,0.497) -- (0.549,0.502) -- (0.615,0.508) -- (0.681,0.515) -- (0.746,0.523) -- (0.810,0.531) -- (0.873,0.541) -- (0.935,0.550) -- (0.996,0.561) -- (1.055,0.572) -- (1.114,0.584) -- (1.171,0.596) -- (1.226,0.609) -- (1.280,0.622) -- (1.333,0.636) -- (1.384,0.651) -- (1.433,0.666) -- (1.480,0.682) -- (1.526,0.698) -- (1.570,0.715) -- (1.611,0.732) -- (1.651,0.750) -- (1.689,0.768) -- (1.725,0.786) -- (1.759,0.805) -- (1.790,0.824) -- (1.820,0.843) -- (1.847,0.863) -- (1.872,0.883) -- (1.894,0.904) -- (1.915,0.924) -- (1.933,0.945) -- (1.948,0.966) -- (1.962,0.987) -- (1.972,1.008) -- (1.981,1.029) -- (1.987,1.051) -- (1.991,1.072) -- (1.992,1.094) -- (1.991,1.115) -- (1.987,1.137) -- (1.981,1.158) -- (1.972,1.179) -- (1.962,1.201);
\draw[grid] (-1.948,1.222) -- (-1.962,1.201) -- (-1.972,1.179) -- (-1.981,1.158) -- (-1.987,1.137) -- (-1.991,1.115) -- (-1.992,1.094) -- (-1.991,1.072) -- (-1.987,1.051) -- (-1.981,1.029) -- (-1.972,1.008) -- (-1.962,0.987) -- (-1.948,0.966) -- (-1.933,0.945) -- (-1.915,0.924) -- (-1.894,0.904) -- (-1.872,0.883) -- (-1.847,0.863) -- (-1.820,0.843) -- (-1.790,0.824) -- (-1.759,0.805) -- (-1.725,0.786) -- (-1.689,0.768) -- (-1.651,0.750) -- (-1.611,0.732) -- (-1.570,0.715) -- (-1.526,0.698) -- (-1.480,0.682) -- (-1.433,0.666) -- (-1.384,0.651) -- (-1.333,0.636) -- (-1.280,0.622) -- (-1.226,0.609) -- (-1.171,0.596) -- (-1.114,0.584) -- (-1.056,0.572) -- (-0.996,0.561);
\draw[grid] (-1.150,-0.615) -- (-1.080,-0.628) -- (-1.008,-0.639) -- (-0.936,-0.649) -- (-0.862,-0.659) -- (-0.787,-0.668) -- (-0.711,-0.676) -- (-0.634,-0.683) -- (-0.556,-0.690) -- (-0.478,-0.695) -- (-0.399,-0.700) -- (-0.320,-0.704) -- (-0.240,-0.707) -- (-0.160,-0.709) -- (-0.080,-0.710) -- (-0.000,-0.711) -- (0.080,-0.710) -- (0.160,-0.709) -- (0.240,-0.707) -- (0.320,-0.704) -- (0.399,-0.700) -- (0.478,-0.695) -- (0.556,-0.690) -- (0.634,-0.683) -- (0.711,-0.676) -- (0.787,-0.668) -- (0.862,-0.659) -- (0.935,-0.649) -- (1.008,-0.639) -- (1.080,-0.628) -- (1.150,-0.615) -- (1.219,-0.603) -- (1.286,-0.589) -- (1.352,-0.575) -- (1.416,-0.560) -- (1.478,-0.544) -- (1.539,-0.528) -- (1.598,-0.511) -- (1.654,-0.494) -- (1.709,-0.476) -- (1.762,-0.457) -- (1.812,-0.438) -- (1.861,-0.418) -- (1.907,-0.397) -- (1.950,-0.377) -- (1.992,-0.355) -- (2.031,-0.334) -- (2.067,-0.312) -- (2.101,-0.289) -- (2.132,-0.266) -- (2.161,-0.243) -- (2.187,-0.220) -- (2.211,-0.196) -- (2.232,-0.172) -- (2.250,-0.148) -- (2.265,-0.123) -- (2.278,-0.099) -- (2.287,-0.074) -- (2.294,-0.050) -- (2.299,-0.025);
\draw[hgrid] (2.299,-0.025) -- (2.300,-0.000) -- (2.299,0.025) -- (2.294,0.050) -- (2.287,0.074) -- (2.278,0.099) -- (2.265,0.123) -- (2.250,0.148) -- (2.232,0.172) -- (2.211,0.196) -- (2.187,0.220) -- (2.161,0.243) -- (2.132,0.266) -- (2.101,0.289) -- (2.067,0.312) -- (2.031,0.334) -- (1.992,0.355) -- (1.950,0.377) -- (1.907,0.397) -- (1.861,0.418) -- (1.812,0.438) -- (1.762,0.457) -- (1.709,0.476) -- (1.654,0.494) -- (1.598,0.511) -- (1.539,0.528) -- (1.478,0.544) -- (1.416,0.560) -- (1.352,0.575) -- (1.286,0.589) -- (1.219,0.603) -- (1.150,0.615) -- (1.080,0.628) -- (1.008,0.639) -- (0.936,0.649) -- (0.862,0.659) -- (0.787,0.668) -- (0.711,0.676) -- (0.634,0.683) -- (0.556,0.690) -- (0.478,0.695) -- (0.399,0.700) -- (0.320,0.704) -- (0.240,0.707) -- (0.160,0.709) -- (0.080,0.710) -- (0.000,0.711) -- (-0.080,0.710) -- (-0.160,0.709) -- (-0.240,0.707) -- (-0.320,0.704) -- (-0.399,0.700) -- (-0.478,0.695) -- (-0.556,0.690) -- (-0.634,0.683) -- (-0.711,0.676) -- (-0.787,0.668) -- (-0.862,0.659) -- (-0.935,0.649) -- (-1.008,0.639) -- (-1.080,0.628) -- (-1.150,0.615) -- (-1.219,0.603) -- (-1.286,0.589) -- (-1.352,0.575) -- (-1.416,0.560) -- (-1.478,0.544) -- (-1.539,0.528) -- (-1.598,0.511) -- (-1.654,0.494) -- (-1.709,0.476) -- (-1.762,0.457) -- (-1.812,0.438) -- (-1.861,0.418) -- (-1.907,0.397) -- (-1.950,0.377) -- (-1.992,0.355) -- (-2.031,0.334) -- (-2.067,0.312) -- (-2.101,0.289) -- (-2.132,0.266) -- (-2.161,0.243) -- (-2.187,0.220) -- (-2.211,0.196) -- (-2.232,0.172) -- (-2.250,0.148) -- (-2.265,0.123) -- (-2.278,0.099) -- (-2.287,0.074) -- (-2.294,0.050) -- (-2.299,0.025);
\draw[grid] (-2.299,0.025) -- (-2.300,0.000) -- (-2.299,-0.025) -- (-2.294,-0.050) -- (-2.287,-0.074) -- (-2.278,-0.099) -- (-2.265,-0.123) -- (-2.250,-0.148) -- (-2.232,-0.172) -- (-2.211,-0.196) -- (-2.187,-0.220) -- (-2.161,-0.243) -- (-2.132,-0.266) -- (-2.101,-0.289) -- (-2.067,-0.312) -- (-2.031,-0.334) -- (-1.992,-0.355) -- (-1.950,-0.377) -- (-1.907,-0.397) -- (-1.861,-0.418) -- (-1.812,-0.438) -- (-1.762,-0.457) -- (-1.709,-0.476) -- (-1.654,-0.494) -- (-1.598,-0.511) -- (-1.539,-0.528) -- (-1.478,-0.544) -- (-1.416,-0.560) -- (-1.352,-0.575) -- (-1.286,-0.589) -- (-1.219,-0.603) -- (-1.150,-0.615);
\draw[grid] (-0.996,-1.627) -- (-0.935,-1.637) -- (-0.873,-1.647) -- (-0.810,-1.656) -- (-0.746,-1.664) -- (-0.681,-1.672) -- (-0.616,-1.679) -- (-0.549,-1.685) -- (-0.482,-1.691) -- (-0.414,-1.696) -- (-0.346,-1.700) -- (-0.277,-1.703) -- (-0.208,-1.706) -- (-0.139,-1.708) -- (-0.070,-1.709) -- (-0.000,-1.709) -- (0.069,-1.709) -- (0.139,-1.708) -- (0.208,-1.706) -- (0.277,-1.703) -- (0.346,-1.700) -- (0.414,-1.696) -- (0.482,-1.691) -- (0.549,-1.685) -- (0.615,-1.679) -- (0.681,-1.672) -- (0.746,-1.664) -- (0.810,-1.656) -- (0.873,-1.647) -- (0.935,-1.637) -- (0.996,-1.627) -- (1.055,-1.616) -- (1.114,-1.604) -- (1.171,-1.592) -- (1.226,-1.579) -- (1.280,-1.565) -- (1.333,-1.551) -- (1.384,-1.537) -- (1.433,-1.521) -- (1.480,-1.506) -- (1.526,-1.489) -- (1.570,-1.473) -- (1.611,-1.456) -- (1.651,-1.438) -- (1.689,-1.420) -- (1.725,-1.402) -- (1.759,-1.383) -- (1.790,-1.364) -- (1.820,-1.344) -- (1.847,-1.324) -- (1.872,-1.304) -- (1.894,-1.284) -- (1.915,-1.263) -- (1.933,-1.243) -- (1.948,-1.222);
\draw[grid] (-1.962,-1.201) -- (-1.948,-1.222) -- (-1.933,-1.243) -- (-1.915,-1.263) -- (-1.894,-1.284) -- (-1.872,-1.304) -- (-1.847,-1.324) -- (-1.820,-1.344) -- (-1.790,-1.364) -- (-1.759,-1.383) -- (-1.725,-1.402) -- (-1.689,-1.420) -- (-1.651,-1.438) -- (-1.611,-1.456) -- (-1.570,-1.473) -- (-1.526,-1.489) -- (-1.480,-1.506) -- (-1.433,-1.521) -- (-1.384,-1.537) -- (-1.333,-1.551) -- (-1.280,-1.565) -- (-1.226,-1.579) -- (-1.171,-1.592) -- (-1.114,-1.604) -- (-1.056,-1.616) -- (-0.996,-1.627);
\draw[grid] (-0.575,-2.202) -- (-0.540,-2.208) -- (-0.504,-2.214) -- (-0.468,-2.219) -- (-0.431,-2.224) -- (-0.393,-2.228) -- (-0.355,-2.232) -- (-0.317,-2.236) -- (-0.278,-2.239) -- (-0.239,-2.242) -- (-0.200,-2.244) -- (-0.160,-2.246) -- (-0.120,-2.248) -- (-0.080,-2.249) -- (-0.040,-2.250) -- (-0.000,-2.250) -- (0.040,-2.250) -- (0.080,-2.249) -- (0.120,-2.248) -- (0.160,-2.246) -- (0.200,-2.244) -- (0.239,-2.242) -- (0.278,-2.239) -- (0.317,-2.236) -- (0.355,-2.232) -- (0.393,-2.228) -- (0.431,-2.224) -- (0.468,-2.219) -- (0.504,-2.214) -- (0.540,-2.208) -- (0.575,-2.202) -- (0.609,-2.196) -- (0.643,-2.189) -- (0.676,-2.182) -- (0.708,-2.174) -- (0.739,-2.167) -- (0.769,-2.159) -- (0.799,-2.150) -- (0.827,-2.141) -- (0.855,-2.132) -- (0.881,-2.123) -- (0.906,-2.113) -- (0.930,-2.103);
\draw[grid] (-0.953,-2.093) -- (-0.930,-2.103) -- (-0.906,-2.113) -- (-0.881,-2.123) -- (-0.855,-2.132) -- (-0.827,-2.141) -- (-0.799,-2.150) -- (-0.769,-2.159) -- (-0.739,-2.167) -- (-0.708,-2.174) -- (-0.676,-2.182) -- (-0.643,-2.189) -- (-0.609,-2.196) -- (-0.575,-2.202);
\draw[grid] (0.000,2.188) -- (-0.020,2.176) -- (-0.040,2.165) -- (-0.060,2.152) -- (-0.080,2.139) -- (-0.100,2.126) -- (-0.120,2.111) -- (-0.140,2.096) -- (-0.160,2.081) -- (-0.180,2.064) -- (-0.200,2.047) -- (-0.219,2.030) -- (-0.239,2.012) -- (-0.259,1.993) -- (-0.278,1.974) -- (-0.298,1.954) -- (-0.317,1.933) -- (-0.336,1.912) -- (-0.355,1.890) -- (-0.374,1.868) -- (-0.393,1.845) -- (-0.412,1.822) -- (-0.431,1.798) -- (-0.449,1.773) -- (-0.468,1.748) -- (-0.486,1.722) -- (-0.504,1.696) -- (-0.522,1.670) -- (-0.540,1.643) -- (-0.558,1.615) -- (-0.575,1.587) -- (-0.592,1.558) -- (-0.609,1.529) -- (-0.626,1.499) -- (-0.643,1.469) -- (-0.660,1.439) -- (-0.676,1.408) -- (-0.692,1.377) -- (-0.708,1.345) -- (-0.724,1.313) -- (-0.739,1.280) -- (-0.754,1.247) -- (-0.770,1.214) -- (-0.784,1.180) -- (-0.799,1.146) -- (-0.813,1.112) -- (-0.827,1.077) -- (-0.841,1.042) -- (-0.855,1.006) -- (-0.868,0.971) -- (-0.881,0.935) -- (-0.894,0.898) -- (-0.906,0.862) -- (-0.918,0.825) -- (-0.930,0.788) -- (-0.942,0.751) -- (-0.953,0.713) -- (-0.964,0.675) -- (-0.975,0.637) -- (-0.986,0.599) -- (-0.996,0.561) -- (-1.006,0.522) -- (-1.015,0.484) -- (-1.025,0.445) -- (-1.034,0.406) -- (-1.042,0.367) -- (-1.051,0.327) -- (-1.059,0.288) -- (-1.066,0.249) -- (-1.074,0.209) -- (-1.081,0.170) -- (-1.087,0.130) -- (-1.094,0.091) -- (-1.100,0.051) -- (-1.105,0.011) -- (-1.111,-0.028) -- (-1.116,-0.068) -- (-1.121,-0.108) -- (-1.125,-0.147) -- (-1.129,-0.187) -- (-1.133,-0.226) -- (-1.136,-0.266) -- (-1.139,-0.305) -- (-1.141,-0.344) -- (-1.144,-0.383) -- (-1.146,-0.422) -- (-1.147,-0.461) -- (-1.148,-0.500) -- (-1.149,-0.539) -- (-1.150,-0.577) -- (-1.150,-0.615) -- (-1.150,-0.654) -- (-1.149,-0.691) -- (-1.148,-0.729) -- (-1.147,-0.767) -- (-1.146,-0.804) -- (-1.144,-0.841) -- (-1.141,-0.877) -- (-1.139,-0.914) -- (-1.136,-0.950) -- (-1.133,-0.986) -- (-1.129,-1.022) -- (-1.125,-1.057) -- (-1.121,-1.092) -- (-1.116,-1.126) -- (-1.111,-1.161) -- (-1.105,-1.195) -- (-1.100,-1.228) -- (-1.094,-1.261) -- (-1.087,-1.294) -- (-1.081,-1.327) -- (-1.074,-1.359) -- (-1.066,-1.390) -- (-1.059,-1.421) -- (-1.051,-1.452) -- (-1.042,-1.482) -- (-1.034,-1.512) -- (-1.025,-1.542) -- (-1.015,-1.570) -- (-1.006,-1.599) -- (-0.996,-1.627) -- (-0.986,-1.654) -- (-0.975,-1.681) -- (-0.964,-1.708) -- (-0.953,-1.734) -- (-0.942,-1.759) -- (-0.930,-1.784) -- (-0.918,-1.808) -- (-0.906,-1.832) -- (-0.894,-1.855) -- (-0.881,-1.878) -- (-0.868,-1.900) -- (-0.855,-1.921) -- (-0.841,-1.942) -- (-0.827,-1.962) -- (-0.813,-1.982) -- (-0.799,-2.001) -- (-0.784,-2.020) -- (-0.770,-2.037) -- (-0.754,-2.055) -- (-0.739,-2.071) -- (-0.724,-2.087) -- (-0.708,-2.103) -- (-0.692,-2.117) -- (-0.676,-2.132) -- (-0.660,-2.145) -- (-0.643,-2.158) -- (-0.626,-2.170) -- (-0.609,-2.181) -- (-0.592,-2.192) -- (-0.575,-2.202) -- (-0.558,-2.212) -- (-0.540,-2.220) -- (-0.522,-2.229) -- (-0.504,-2.236) -- (-0.486,-2.243) -- (-0.468,-2.249) -- (-0.449,-2.254) -- (-0.431,-2.259) -- (-0.412,-2.263);
\draw[grid] (0.000,2.188) -- (-0.000,2.175) -- (-0.000,2.161) -- (-0.000,2.147) -- (-0.000,2.133) -- (-0.000,2.117) -- (-0.000,2.101) -- (-0.000,2.085) -- (-0.000,2.067) -- (-0.000,2.049) -- (-0.000,2.031) -- (-0.000,2.012) -- (-0.000,1.992) -- (-0.000,1.972) -- (-0.000,1.951) -- (-0.000,1.929) -- (-0.000,1.907) -- (-0.000,1.884) -- (-0.000,1.861) -- (-0.000,1.837) -- (-0.000,1.813) -- (-0.000,1.788) -- (-0.000,1.762) -- (-0.000,1.736) -- (-0.000,1.709) -- (-0.000,1.682) -- (-0.000,1.655) -- (-0.000,1.626) -- (-0.000,1.598) -- (-0.000,1.569) -- (-0.000,1.539) -- (-0.000,1.509) -- (-0.000,1.479) -- (-0.000,1.448) -- (-0.000,1.416) -- (-0.000,1.384) -- (-0.000,1.352) -- (-0.000,1.319) -- (-0.000,1.286) -- (-0.000,1.253) -- (-0.000,1.219) -- (-0.000,1.185) -- (-0.000,1.150) -- (-0.000,1.115) -- (-0.000,1.080) -- (-0.000,1.044) -- (-0.000,1.008) -- (-0.000,0.972) -- (-0.000,0.936) -- (-0.000,0.899) -- (-0.000,0.862) -- (-0.000,0.824) -- (-0.000,0.787) -- (-0.000,0.749) -- (-0.000,0.711) -- (-0.000,0.673) -- (-0.000,0.634) -- (-0.000,0.595) -- (-0.000,0.557) -- (-0.000,0.517) -- (-0.000,0.478) -- (-0.000,0.439) -- (-0.000,0.399) -- (-0.000,0.360) -- (-0.000,0.320) -- (-0.000,0.280) -- (-0.000,0.240) -- (-0.000,0.201) -- (-0.000,0.161) -- (-0.000,0.120) -- (-0.000,0.080) -- (-0.000,0.040) -- (-0.000,0.000) -- (-0.000,-0.040) -- (-0.000,-0.080) -- (-0.000,-0.120) -- (-0.000,-0.160) -- (-0.000,-0.200) -- (-0.000,-0.240) -- (-0.000,-0.280) -- (-0.000,-0.320) -- (-0.000,-0.360) -- (-0.000,-0.399) -- (-0.000,-0.439) -- (-0.000,-0.478) -- (-0.000,-0.517) -- (-0.000,-0.556) -- (-0.000,-0.595) -- (-0.000,-0.634) -- (-0.000,-0.672) -- (-0.000,-0.711) -- (-0.000,-0.749) -- (-0.000,-0.787) -- (-0.000,-0.824) -- (-0.000,-0.862) -- (-0.000,-0.899) -- (-0.000,-0.935) -- (-0.000,-0.972) -- (-0.000,-1.008) -- (-0.000,-1.044) -- (-0.000,-1.080) -- (-0.000,-1.115) -- (-0.000,-1.150) -- (-0.000,-1.185) -- (-0.000,-1.219) -- (-0.000,-1.253) -- (-0.000,-1.286) -- (-0.000,-1.319) -- (-0.000,-1.352) -- (-0.000,-1.384) -- (-0.000,-1.416) -- (-0.000,-1.447) -- (-0.000,-1.478) -- (-0.000,-1.509) -- (-0.000,-1.539) -- (-0.000,-1.569) -- (-0.000,-1.598) -- (-0.000,-1.626) -- (-0.000,-1.654) -- (-0.000,-1.682) -- (-0.000,-1.709) -- (-0.000,-1.736) -- (-0.000,-1.762) -- (-0.000,-1.787) -- (-0.000,-1.812) -- (-0.000,-1.837) -- (-0.000,-1.861) -- (-0.000,-1.884) -- (-0.000,-1.907) -- (-0.000,-1.929) -- (-0.000,-1.951) -- (-0.000,-1.972) -- (-0.000,-1.992) -- (-0.000,-2.012) -- (-0.000,-2.031) -- (-0.000,-2.049) -- (-0.000,-2.067) -- (-0.000,-2.085) -- (-0.000,-2.101) -- (-0.000,-2.117) -- (-0.000,-2.133) -- (-0.000,-2.147) -- (-0.000,-2.161) -- (-0.000,-2.175) -- (-0.000,-2.187) -- (-0.000,-2.200) -- (-0.000,-2.211) -- (-0.000,-2.222) -- (-0.000,-2.232) -- (-0.000,-2.241) -- (-0.000,-2.250) -- (-0.000,-2.258) -- (-0.000,-2.265) -- (-0.000,-2.272) -- (-0.000,-2.278) -- (-0.000,-2.283) -- (-0.000,-2.287) -- (-0.000,-2.291) -- (-0.000,-2.294) -- (-0.000,-2.297) -- (-0.000,-2.299) -- (-0.000,-2.300) -- (-0.000,-2.300);
\draw[grid] (0.000,2.188) -- (0.020,2.176) -- (0.040,2.165) -- (0.060,2.152) -- (0.080,2.139) -- (0.100,2.126) -- (0.120,2.111) -- (0.140,2.096) -- (0.160,2.081) -- (0.180,2.064) -- (0.200,2.047) -- (0.219,2.030) -- (0.239,2.012) -- (0.259,1.993) -- (0.278,1.974) -- (0.298,1.954) -- (0.317,1.933) -- (0.336,1.912) -- (0.355,1.890) -- (0.374,1.868) -- (0.393,1.845) -- (0.412,1.822) -- (0.431,1.798) -- (0.449,1.773) -- (0.468,1.748) -- (0.486,1.722) -- (0.504,1.696) -- (0.522,1.670) -- (0.540,1.643) -- (0.558,1.615) -- (0.575,1.587) -- (0.592,1.558) -- (0.609,1.529) -- (0.626,1.499) -- (0.643,1.469) -- (0.660,1.439) -- (0.676,1.408) -- (0.692,1.377) -- (0.708,1.345) -- (0.724,1.313) -- (0.739,1.280) -- (0.754,1.247) -- (0.769,1.214) -- (0.784,1.180) -- (0.799,1.146) -- (0.813,1.112) -- (0.827,1.077) -- (0.841,1.042) -- (0.855,1.006) -- (0.868,0.971) -- (0.881,0.935) -- (0.894,0.898) -- (0.906,0.862) -- (0.918,0.825) -- (0.930,0.788) -- (0.942,0.751) -- (0.953,0.713) -- (0.964,0.675) -- (0.975,0.637) -- (0.986,0.599) -- (0.996,0.561) -- (1.006,0.522) -- (1.015,0.484) -- (1.025,0.445) -- (1.034,0.406) -- (1.042,0.367) -- (1.051,0.327) -- (1.059,0.288) -- (1.066,0.249) -- (1.074,0.209) -- (1.081,0.170) -- (1.087,0.130) -- (1.094,0.091) -- (1.100,0.051) -- (1.105,0.011) -- (1.111,-0.028) -- (1.116,-0.068) -- (1.120,-0.108) -- (1.125,-0.147) -- (1.129,-0.187) -- (1.132,-0.226) -- (1.136,-0.266) -- (1.139,-0.305) -- (1.141,-0.344) -- (1.144,-0.383) -- (1.146,-0.422) -- (1.147,-0.461) -- (1.148,-0.500) -- (1.149,-0.539) -- (1.150,-0.577) -- (1.150,-0.615) -- (1.150,-0.654) -- (1.149,-0.691) -- (1.148,-0.729) -- (1.147,-0.767) -- (1.146,-0.804) -- (1.144,-0.841) -- (1.141,-0.877) -- (1.139,-0.914) -- (1.136,-0.950) -- (1.132,-0.986) -- (1.129,-1.022) -- (1.125,-1.057) -- (1.120,-1.092) -- (1.116,-1.126) -- (1.111,-1.161) -- (1.105,-1.195) -- (1.100,-1.228) -- (1.094,-1.261) -- (1.087,-1.294) -- (1.081,-1.327) -- (1.074,-1.359) -- (1.066,-1.390) -- (1.059,-1.421) -- (1.051,-1.452) -- (1.042,-1.482) -- (1.034,-1.512) -- (1.025,-1.542) -- (1.015,-1.570) -- (1.006,-1.599) -- (0.996,-1.627) -- (0.986,-1.654) -- (0.975,-1.681) -- (0.964,-1.708) -- (0.953,-1.734) -- (0.942,-1.759) -- (0.930,-1.784) -- (0.918,-1.808) -- (0.906,-1.832) -- (0.894,-1.855) -- (0.881,-1.878) -- (0.868,-1.900) -- (0.855,-1.921) -- (0.841,-1.942) -- (0.827,-1.962) -- (0.813,-1.982) -- (0.799,-2.001) -- (0.784,-2.020) -- (0.769,-2.037) -- (0.754,-2.055) -- (0.739,-2.071) -- (0.724,-2.087) -- (0.708,-2.103) -- (0.692,-2.117) -- (0.676,-2.132) -- (0.660,-2.145) -- (0.643,-2.158) -- (0.626,-2.170) -- (0.609,-2.181) -- (0.592,-2.192) -- (0.575,-2.202) -- (0.558,-2.212) -- (0.540,-2.220) -- (0.522,-2.229) -- (0.504,-2.236) -- (0.486,-2.243) -- (0.468,-2.249) -- (0.449,-2.254) -- (0.431,-2.259) -- (0.412,-2.263);
\draw[grid] (0.000,2.188) -- (0.035,2.181) -- (0.070,2.174) -- (0.104,2.166) -- (0.139,2.157) -- (0.174,2.148) -- (0.208,2.138) -- (0.243,2.128) -- (0.277,2.117) -- (0.312,2.105) -- (0.346,2.093) -- (0.380,2.080) -- (0.414,2.066) -- (0.448,2.052) -- (0.482,2.037) -- (0.516,2.021) -- (0.549,2.005) -- (0.582,1.988) -- (0.616,1.971) -- (0.648,1.953) -- (0.681,1.934) -- (0.714,1.915) -- (0.746,1.895) -- (0.778,1.875) -- (0.810,1.854) -- (0.842,1.832) -- (0.873,1.810) -- (0.904,1.788) -- (0.935,1.765) -- (0.966,1.741) -- (0.996,1.717) -- (1.026,1.692) -- (1.055,1.667) -- (1.085,1.641) -- (1.114,1.615) -- (1.142,1.588) -- (1.171,1.561) -- (1.199,1.533) -- (1.226,1.505) -- (1.253,1.476) -- (1.280,1.447) -- (1.307,1.418) -- (1.333,1.388) -- (1.358,1.358) -- (1.384,1.327) -- (1.408,1.296) -- (1.433,1.264) -- (1.457,1.232) -- (1.480,1.200) -- (1.503,1.167) -- (1.526,1.134) -- (1.548,1.100) -- (1.570,1.067) -- (1.591,1.033) -- (1.611,0.998) -- (1.632,0.964) -- (1.651,0.929) -- (1.670,0.893) -- (1.689,0.858) -- (1.707,0.822) -- (1.725,0.786) -- (1.742,0.750) -- (1.759,0.713) -- (1.775,0.676) -- (1.790,0.640) -- (1.805,0.602) -- (1.820,0.565) -- (1.833,0.528) -- (1.847,0.490) -- (1.860,0.452) -- (1.872,0.414) -- (1.883,0.376) -- (1.894,0.338) -- (1.905,0.300) -- (1.915,0.261) -- (1.924,0.223) -- (1.933,0.184) -- (1.941,0.146) -- (1.948,0.107) -- (1.955,0.069) -- (1.962,0.030) -- (1.967,-0.009) -- (1.972,-0.047) -- (1.977,-0.086) -- (1.981,-0.125) -- (1.984,-0.163) -- (1.987,-0.202) -- (1.989,-0.240) -- (1.991,-0.279) -- (1.991,-0.317) -- (1.992,-0.355) -- (1.991,-0.393) -- (1.991,-0.431) -- (1.989,-0.469) -- (1.987,-0.507) -- (1.984,-0.545) -- (1.981,-0.582) -- (1.977,-0.619) -- (1.972,-0.656) -- (1.967,-0.693) -- (1.962,-0.730) -- (1.955,-0.766) -- (1.948,-0.802) -- (1.941,-0.838) -- (1.933,-0.874) -- (1.924,-0.909) -- (1.915,-0.945) -- (1.905,-0.979) -- (1.894,-1.014) -- (1.883,-1.048) -- (1.872,-1.082) -- (1.860,-1.116) -- (1.847,-1.149) -- (1.833,-1.182) -- (1.820,-1.214) -- (1.805,-1.247) -- (1.790,-1.278) -- (1.775,-1.310) -- (1.759,-1.341) -- (1.742,-1.371) -- (1.725,-1.402) -- (1.707,-1.431) -- (1.689,-1.461) -- (1.670,-1.489) -- (1.651,-1.518) -- (1.632,-1.546) -- (1.611,-1.573) -- (1.591,-1.600) -- (1.570,-1.627) -- (1.548,-1.653) -- (1.526,-1.678) -- (1.503,-1.703) -- (1.480,-1.728) -- (1.457,-1.752) -- (1.433,-1.775) -- (1.408,-1.798) -- (1.384,-1.820) -- (1.358,-1.842) -- (1.333,-1.863) -- (1.307,-1.884) -- (1.280,-1.904) -- (1.253,-1.924) -- (1.226,-1.943) -- (1.199,-1.961) -- (1.171,-1.979) -- (1.142,-1.996) -- (1.114,-2.012);
\draw[grid] (0.000,2.188) -- (0.040,2.187) -- (0.080,2.186) -- (0.120,2.185) -- (0.160,2.182) -- (0.200,2.179) -- (0.240,2.176) -- (0.280,2.171) -- (0.320,2.166) -- (0.360,2.161) -- (0.399,2.154) -- (0.439,2.147) -- (0.478,2.140) -- (0.517,2.131) -- (0.556,2.123) -- (0.595,2.113) -- (0.634,2.103) -- (0.672,2.092) -- (0.711,2.080) -- (0.749,2.068) -- (0.787,2.056) -- (0.824,2.042) -- (0.862,2.028) -- (0.899,2.014) -- (0.935,1.998) -- (0.972,1.983) -- (1.008,1.966) -- (1.044,1.949) -- (1.080,1.931) -- (1.115,1.913) -- (1.150,1.894) -- (1.185,1.875) -- (1.219,1.855) -- (1.253,1.835) -- (1.286,1.814) -- (1.319,1.792) -- (1.352,1.770) -- (1.384,1.747) -- (1.416,1.724) -- (1.447,1.700) -- (1.478,1.676) -- (1.509,1.651) -- (1.539,1.626) -- (1.569,1.600) -- (1.598,1.574) -- (1.626,1.547) -- (1.654,1.520) -- (1.682,1.492) -- (1.709,1.464) -- (1.736,1.435) -- (1.762,1.406) -- (1.787,1.377) -- (1.812,1.347) -- (1.837,1.316) -- (1.861,1.286) -- (1.884,1.255) -- (1.907,1.223) -- (1.929,1.191) -- (1.950,1.159) -- (1.971,1.127) -- (1.992,1.094) -- (2.012,1.061) -- (2.031,1.027) -- (2.049,0.993) -- (2.067,0.959) -- (2.084,0.924) -- (2.101,0.890) -- (2.117,0.855) -- (2.132,0.819) -- (2.147,0.784) -- (2.161,0.748) -- (2.175,0.712) -- (2.187,0.676) -- (2.199,0.640) -- (2.211,0.603) -- (2.222,0.566) -- (2.232,0.529) -- (2.241,0.492) -- (2.250,0.455) -- (2.258,0.417) -- (2.265,0.380) -- (2.272,0.342) -- (2.278,0.304) -- (2.283,0.267) -- (2.287,0.229) -- (2.291,0.191) -- (2.294,0.153) -- (2.297,0.114) -- (2.299,0.076) -- (2.300,0.038);
\draw[grid] (0.000,2.188) -- (0.035,2.193) -- (0.070,2.199) -- (0.104,2.203) -- (0.139,2.207) -- (0.174,2.210) -- (0.208,2.213) -- (0.243,2.215) -- (0.277,2.216) -- (0.312,2.216) -- (0.346,2.216) -- (0.380,2.215) -- (0.414,2.214) -- (0.448,2.211) -- (0.482,2.209) -- (0.516,2.205) -- (0.549,2.201) -- (0.582,2.196) -- (0.616,2.190) -- (0.648,2.184) -- (0.681,2.177) -- (0.714,2.170) -- (0.746,2.161) -- (0.778,2.152) -- (0.810,2.143) -- (0.842,2.133) -- (0.873,2.122) -- (0.904,2.110) -- (0.935,2.098) -- (0.966,2.086) -- (0.996,2.072) -- (1.026,2.058) -- (1.056,2.043) -- (1.085,2.028);
\draw[grid] (0.000,2.188) -- (0.020,2.198) -- (0.040,2.208) -- (0.060,2.217) -- (0.080,2.225) -- (0.100,2.233) -- (0.120,2.240) -- (0.140,2.246) -- (0.160,2.252) -- (0.180,2.257) -- (0.200,2.261) -- (0.219,2.265) -- (0.239,2.268) -- (0.259,2.270) -- (0.278,2.271) -- (0.298,2.272) -- (0.317,2.272) -- (0.336,2.272) -- (0.355,2.271) -- (0.374,2.269) -- (0.393,2.266);
\draw[grid] (0.000,2.188) -- (0.000,2.200) -- (0.000,2.211) -- (0.000,2.222) -- (0.000,2.232) -- (0.000,2.241) -- (0.000,2.250) -- (0.000,2.258) -- (0.000,2.265) -- (0.000,2.272) -- (0.000,2.278) -- (0.000,2.283) -- (0.000,2.287) -- (0.000,2.291) -- (0.000,2.294) -- (0.000,2.297) -- (0.000,2.299) -- (0.000,2.300);
\draw[grid] (0.000,2.188) -- (-0.020,2.198) -- (-0.040,2.208) -- (-0.060,2.217) -- (-0.080,2.225) -- (-0.100,2.233) -- (-0.120,2.240) -- (-0.140,2.246) -- (-0.160,2.252) -- (-0.180,2.257) -- (-0.200,2.261) -- (-0.219,2.265) -- (-0.239,2.268) -- (-0.259,2.270) -- (-0.278,2.271) -- (-0.298,2.272) -- (-0.317,2.272) -- (-0.336,2.272) -- (-0.355,2.271) -- (-0.374,2.269) -- (-0.393,2.266);
\draw[grid] (0.000,2.188) -- (-0.035,2.193) -- (-0.070,2.199) -- (-0.104,2.203) -- (-0.139,2.207) -- (-0.174,2.210) -- (-0.208,2.213) -- (-0.243,2.215) -- (-0.277,2.216) -- (-0.312,2.216) -- (-0.346,2.216) -- (-0.380,2.215) -- (-0.414,2.214) -- (-0.448,2.211) -- (-0.482,2.209) -- (-0.516,2.205) -- (-0.549,2.201) -- (-0.582,2.196) -- (-0.616,2.190) -- (-0.648,2.184) -- (-0.681,2.177) -- (-0.714,2.170) -- (-0.746,2.161) -- (-0.778,2.152) -- (-0.810,2.143) -- (-0.842,2.133) -- (-0.873,2.122) -- (-0.904,2.110) -- (-0.935,2.098) -- (-0.966,2.086) -- (-0.996,2.072) -- (-1.026,2.058) -- (-1.055,2.043) -- (-1.085,2.028);
\draw[grid] (0.000,2.188) -- (-0.040,2.187) -- (-0.080,2.186) -- (-0.120,2.185) -- (-0.160,2.182) -- (-0.200,2.179) -- (-0.240,2.176) -- (-0.280,2.171) -- (-0.320,2.166) -- (-0.360,2.161) -- (-0.399,2.154) -- (-0.439,2.147) -- (-0.478,2.140) -- (-0.517,2.131) -- (-0.556,2.123) -- (-0.595,2.113) -- (-0.634,2.103) -- (-0.672,2.092) -- (-0.711,2.080) -- (-0.749,2.068) -- (-0.787,2.056) -- (-0.824,2.042) -- (-0.862,2.028) -- (-0.899,2.014) -- (-0.935,1.998) -- (-0.972,1.983) -- (-1.008,1.966) -- (-1.044,1.949) -- (-1.080,1.931) -- (-1.115,1.913) -- (-1.150,1.894) -- (-1.185,1.875) -- (-1.219,1.855) -- (-1.253,1.835) -- (-1.286,1.814) -- (-1.319,1.792) -- (-1.352,1.770) -- (-1.384,1.747) -- (-1.416,1.724) -- (-1.447,1.700) -- (-1.478,1.676) -- (-1.509,1.651) -- (-1.539,1.626) -- (-1.569,1.600) -- (-1.598,1.574) -- (-1.626,1.547) -- (-1.654,1.520) -- (-1.682,1.492) -- (-1.709,1.464) -- (-1.736,1.435) -- (-1.762,1.406) -- (-1.787,1.377) -- (-1.812,1.347) -- (-1.837,1.316) -- (-1.861,1.286) -- (-1.884,1.255) -- (-1.907,1.223) -- (-1.929,1.191) -- (-1.950,1.159) -- (-1.971,1.127) -- (-1.992,1.094) -- (-2.012,1.061) -- (-2.031,1.027) -- (-2.049,0.993) -- (-2.067,0.959) -- (-2.084,0.924) -- (-2.101,0.890) -- (-2.117,0.855) -- (-2.132,0.819) -- (-2.147,0.784) -- (-2.161,0.748) -- (-2.175,0.712) -- (-2.187,0.676) -- (-2.199,0.640) -- (-2.211,0.603) -- (-2.222,0.566) -- (-2.232,0.529) -- (-2.241,0.492) -- (-2.250,0.455) -- (-2.258,0.417) -- (-2.265,0.380) -- (-2.272,0.342) -- (-2.278,0.304) -- (-2.283,0.267) -- (-2.287,0.229) -- (-2.291,0.191) -- (-2.294,0.153) -- (-2.297,0.114) -- (-2.299,0.076) -- (-2.300,0.038) -- (-2.300,0.000);
\draw[grid] (0.000,2.188) -- (-0.035,2.181) -- (-0.070,2.174) -- (-0.104,2.166) -- (-0.139,2.157) -- (-0.174,2.148) -- (-0.208,2.138) -- (-0.243,2.128) -- (-0.277,2.117) -- (-0.312,2.105) -- (-0.346,2.093) -- (-0.380,2.080) -- (-0.414,2.066) -- (-0.448,2.052) -- (-0.482,2.037) -- (-0.516,2.021) -- (-0.549,2.005) -- (-0.582,1.988) -- (-0.616,1.971) -- (-0.648,1.953) -- (-0.681,1.934) -- (-0.714,1.915) -- (-0.746,1.895) -- (-0.778,1.875) -- (-0.810,1.854) -- (-0.842,1.832) -- (-0.873,1.810) -- (-0.904,1.788) -- (-0.935,1.765) -- (-0.966,1.741) -- (-0.996,1.717) -- (-1.026,1.692) -- (-1.056,1.667) -- (-1.085,1.641) -- (-1.114,1.615) -- (-1.142,1.588) -- (-1.171,1.561) -- (-1.199,1.533) -- (-1.226,1.505) -- (-1.253,1.476) -- (-1.280,1.447) -- (-1.307,1.418) -- (-1.333,1.388) -- (-1.358,1.358) -- (-1.384,1.327) -- (-1.408,1.296) -- (-1.433,1.264) -- (-1.457,1.232) -- (-1.480,1.200) -- (-1.503,1.167) -- (-1.526,1.134) -- (-1.548,1.101) -- (-1.570,1.067) -- (-1.591,1.033) -- (-1.611,0.998) -- (-1.632,0.964) -- (-1.651,0.929) -- (-1.670,0.893) -- (-1.689,0.858) -- (-1.707,0.822) -- (-1.725,0.786) -- (-1.742,0.750) -- (-1.759,0.713) -- (-1.775,0.677) -- (-1.790,0.640) -- (-1.805,0.602) -- (-1.820,0.565) -- (-1.833,0.528) -- (-1.847,0.490) -- (-1.860,0.452) -- (-1.872,0.414) -- (-1.883,0.376) -- (-1.894,0.338) -- (-1.905,0.300) -- (-1.915,0.261) -- (-1.924,0.223) -- (-1.933,0.184) -- (-1.941,0.146) -- (-1.948,0.107) -- (-1.955,0.069) -- (-1.962,0.030) -- (-1.967,-0.009) -- (-1.972,-0.047) -- (-1.977,-0.086) -- (-1.981,-0.125) -- (-1.984,-0.163) -- (-1.987,-0.202) -- (-1.989,-0.240) -- (-1.991,-0.279) -- (-1.992,-0.317) -- (-1.992,-0.355) -- (-1.992,-0.393) -- (-1.991,-0.431) -- (-1.989,-0.469) -- (-1.987,-0.507) -- (-1.984,-0.545) -- (-1.981,-0.582) -- (-1.977,-0.619) -- (-1.972,-0.656) -- (-1.967,-0.693) -- (-1.962,-0.730) -- (-1.955,-0.766) -- (-1.948,-0.802) -- (-1.941,-0.838) -- (-1.933,-0.874) -- (-1.924,-0.909) -- (-1.915,-0.945) -- (-1.905,-0.979) -- (-1.894,-1.014) -- (-1.883,-1.048) -- (-1.872,-1.082) -- (-1.860,-1.116) -- (-1.847,-1.149) -- (-1.833,-1.182) -- (-1.820,-1.214) -- (-1.805,-1.247) -- (-1.790,-1.278) -- (-1.775,-1.310) -- (-1.759,-1.341) -- (-1.742,-1.371) -- (-1.725,-1.402) -- (-1.707,-1.431) -- (-1.689,-1.461) -- (-1.670,-1.489) -- (-1.651,-1.518) -- (-1.632,-1.546) -- (-1.611,-1.573) -- (-1.591,-1.600) -- (-1.570,-1.627) -- (-1.548,-1.653) -- (-1.526,-1.678) -- (-1.503,-1.703) -- (-1.480,-1.728) -- (-1.457,-1.752) -- (-1.433,-1.775) -- (-1.408,-1.798) -- (-1.384,-1.820) -- (-1.358,-1.842) -- (-1.333,-1.863) -- (-1.307,-1.884) -- (-1.280,-1.904) -- (-1.253,-1.924) -- (-1.226,-1.943) -- (-1.199,-1.961) -- (-1.171,-1.979) -- (-1.142,-1.996) -- (-1.114,-2.012);
\draw[celldiag] (-0.000,-0.711) -- (0.060,-0.767) -- (0.120,-0.823) -- (0.180,-0.878) -- (0.239,-0.932) -- (0.298,-0.984) -- (0.355,-1.036) -- (0.412,-1.086) -- (0.468,-1.135) -- (0.522,-1.183) -- (0.575,-1.229) -- (0.626,-1.275) -- (0.676,-1.319) -- (0.724,-1.362) -- (0.769,-1.403) -- (0.813,-1.444) -- (0.855,-1.483) -- (0.894,-1.521) -- (0.930,-1.557) -- (0.964,-1.593) -- (0.996,-1.627);
\fill[cellpt] (-0.000,-0.711) circle (1.1pt) node[anchor=north east,inner sep=1pt,font=\scriptsize\itshape] {a};
\fill[cellpt] (1.150,-0.615) circle (1.1pt) node[anchor=north west,inner sep=1pt,font=\scriptsize\itshape] {b};
\fill[cellpt] (-0.000,-1.709) circle (1.1pt) node[anchor=south east,inner sep=1pt,font=\scriptsize\itshape] {c};
\fill[cellpt] (0.996,-1.627) circle (1.1pt) node[anchor=south west,inner sep=1pt,font=\scriptsize\itshape] {d};
\draw[axisin] (0,0) -- (1.992,-0.355);
\draw[axis] (1.992,-0.355) -- (2.888,-0.515) node[anchor=north west] {$x$};
\draw[axisin] (0,0) -- (0.000,2.188);
\draw[axis] (0.000,2.188) -- (0.000,3.172) node[anchor=south] {$y$};
\draw[axisin] (0,0) -- (-1.150,-0.615);
\draw[axis] (-1.150,-0.615) -- (-1.667,-0.892) node[anchor=north east] {$z$};
\draw[helper] (0.640,0.974) -- (0.640,-0.490);
\draw[helper] (0,0) -- (0.640,-0.490) ;
\draw[radius] (0,0) -- node[pos=0.3, below right=-1pt, font=\small] {$r$} (0.640,0.974);
\draw[thetaarc,->] (0.000,1.203) -- (0.020,1.187) -- (0.040,1.169) -- (0.059,1.148) -- (0.079,1.126) -- (0.099,1.101) -- (0.118,1.075) -- (0.137,1.047) -- (0.156,1.017) -- (0.174,0.985) -- (0.193,0.952) -- (0.211,0.917) -- (0.228,0.880) -- (0.245,0.841) -- (0.262,0.802) -- (0.279,0.760) -- (0.294,0.718) -- (0.310,0.674) -- (0.324,0.629) -- (0.339,0.583) -- (0.352,0.536);
\node[thetacol] at (0.252,1.246) {$\theta$};
\draw[phiarc,->] (-0.575,-0.308) -- (-0.529,-0.315) -- (-0.482,-0.323) -- (-0.435,-0.329) -- (-0.386,-0.335) -- (-0.336,-0.340) -- (-0.286,-0.344) -- (-0.235,-0.348) -- (-0.184,-0.351) -- (-0.132,-0.353) -- (-0.080,-0.354) -- (-0.028,-0.355) -- (0.024,-0.355) -- (0.076,-0.355) -- (0.128,-0.353) -- (0.180,-0.351) -- (0.231,-0.348) -- (0.282,-0.344) -- (0.332,-0.340) -- (0.382,-0.335) -- (0.431,-0.329);
\node[phicol] at (-0.106,-0.468) {$\varphi$};
\node[font=\scriptsize, anchor=west] at (0.690,0.242) {$r\cos\theta$};
\node[font=\scriptsize, anchor=north west] at (0.384,-0.314) {$r\sin\theta$};
\fill[black] (0.640,0.974) circle (1.8pt) node[anchor=south west, inner sep=2pt] {$P$};
\fill[black] (0,0) circle (1.2pt) node[anchor=north east, inner sep=1.5pt, font=\small] {$O$};
\fill[polecol] (0.000,2.188) circle (1.6pt);
\node[font=\scriptsize, anchor=south east] at (0.000,2.188) {north pole $\theta{=}0$};
\fill[polecol] (-0.000,-2.188) circle (1.6pt) node[anchor=north, font=\scriptsize] {south pole $\theta{=}\pi$};
\draw[line width=0.8pt] (0,0) circle (2.3);
\end{tikzpicture}\par\vspace{2pt}
{\small\textbf{Figure 1.} Any point $P$ on the sphere is picked by two angles:
{\color{thetac}$\theta$} (down from the north pole) and {\color{phic}$\varphi$} (around the $y$ axis, starting at $+z$).
The {\color{cellc!60!black}orange} patch is one grid cell.}
\end{minipage}\hfill
\begin{minipage}[c]{0.53\textwidth}\centering
\begin{tikzpicture}[x=0.78cm,y=0.78cm, font=\scriptsize, alt={Grid diagram of the sphere unwrapped into angle space, drawn with 4 rows and 8 columns of cells. Columns run left to right in phi from 0 to 2 pi (column j = 0 to LON_BANDS); rows run top to bottom in theta, row i = 0 to 4, from the north pole at the top to the south pole at the bottom. Each cell is split by a diagonal into two triangles. In the top row, the triangle touching the north pole edge is hatched as skipped because it has zero area, and the other triangle is drawn. In the bottom row, the triangle touching the south pole edge is hatched as skipped. Both triangles are drawn in the middle rows. The left and right edges are highlighted in blue as the seam, where the columns phi = 0 and phi = 2 pi are the same points. The cell in row 2, column 2 is shaded orange with corners a (top left), b (top right), c (bottom left), and d (bottom right), and curved arrows show that both triangles a, c, d and a, d, b are wound counter-clockwise.}]
  \def\LAT{4}\def\LON{8}
  % pole rows: one kept triangle, one skipped triangle
  \foreach \j in {0,...,7} {
    % top row (i = 0): a == b, keep (a,c,d), skip (a,d,b)
    \fill[keep!25] (\j,0) -- (\j,-1) -- (\j+1,-1) -- cycle;
    \fill[pattern=north east lines, pattern color=gray!60] (\j,0) -- (\j+1,-1) -- (\j+1,0) -- cycle;
    % bottom row (i = 3): c == d, keep (a,d,b), skip (a,c,d)
    \fill[keep!25] (\j,-3) -- (\j+1,-4) -- (\j+1,-3) -- cycle;
    \fill[pattern=north east lines, pattern color=gray!60] (\j,-3) -- (\j,-4) -- (\j+1,-4) -- cycle;
  }
  % middle rows: both triangles kept
  \fill[keep!12] (0,-1) rectangle (8,-3);
  % highlighted cell i = 2, j = 2
  \fill[cellc!30] (2,-2) rectangle (3,-3);
  % grid lines and diagonals
  \foreach \j in {1,...,7} \draw[gray!70] (\j,0) -- (\j,-4);
  \foreach \i in {1,...,3} \draw[gray!70] (0,-\i) -- (8,-\i);
  \foreach \i in {0,...,3} \foreach \j in {0,...,7} \draw[gray!70] (\j,-\i) -- (\j+1,-\i-1);
  \draw[line width=0.8pt] (0,0) rectangle (8,-4);
  % seam column
  \draw[phic, line width=1.2pt] (0,0) -- (0,-4);
  \draw[phic, line width=1.2pt] (8,0) -- (8,-4);
  % vertices
  \foreach \i in {0,...,4} \foreach \j in {0,...,8} \fill (\j,-\i) circle (1pt);
  % highlighted cell labels
  \draw[cellc!80!black, line width=1pt] (2,-2) rectangle (3,-3);
  \draw[cellc!80!black, line width=1pt] (2,-2) -- (3,-3);
  \node[anchor=south east, inner sep=1pt, font=\scriptsize\itshape\bfseries] at (2,-2) {a};
  \node[anchor=south west, inner sep=1pt, font=\scriptsize\itshape\bfseries] at (3,-2) {b};
  \node[anchor=north east, inner sep=1pt, font=\scriptsize\itshape\bfseries] at (2,-3) {c};
  \node[anchor=north west, inner sep=1pt, font=\scriptsize\itshape\bfseries] at (3,-3) {d};
  \draw[-{Stealth[length=3pt]}, black!80] ($(2,-2)+(0.33,-0.52)$) arc[start angle=-150, end angle=160, radius=0.16];
  \draw[-{Stealth[length=3pt]}, black!80] ($(2,-2)+(0.80,-0.30)$) arc[start angle=-150, end angle=160, radius=0.16];
  % axes labels
  \node[anchor=south] at (0,0.05) {$0$};
  \node[anchor=south] at (4,0.05) {$\pi$};
  \node[anchor=south] at (8,0.05) {$2\pi$};
  \draw[-{Stealth[length=4pt]}, phic, line width=0.8pt] (0.6,0.55) -- (7.4,0.55) node[midway, above, font=\small] {$\varphi$ \;(column $j = 0 \dots \text{LON\_BANDS}$)};
  \foreach \i/\t in {0/0, 1/\frac{\pi}{4}, 2/\frac{\pi}{2}, 3/\frac{3\pi}{4}, 4/\pi} \node[anchor=east] at (-0.08,-\i) {$\i$};
  \draw[-{Stealth[length=4pt]}, thetac, line width=0.8pt] (-0.75,-0.3) -- (-0.75,-3.7) node[midway, above, rotate=90, font=\small] {$\theta$ \;(row $i$)};
  \node[anchor=west, align=left] at (8.15,0) {north pole};
  \node[anchor=west, align=left] at (8.15,-4) {south pole};
  % legend
  \begin{scope}[shift={(0,-4.75)}]
    \fill[keep!25] (0,0) rectangle (0.4,-0.3); \draw[gray!70] (0,0) rectangle (0.4,-0.3);
    \node[anchor=west] at (0.45,-0.15) {drawn};
    \fill[pattern=north east lines, pattern color=gray!60] (1.9,0) rectangle (2.3,-0.3); \draw[gray!70] (1.9,0) rectangle (2.3,-0.3);
    \node[anchor=west] at (2.35,-0.15) {skipped (zero area)};
    \draw[phic, line width=1.2pt] (6.0,0) -- (6.0,-0.3);
    \node[anchor=west] at (6.05,-0.15) {seam: same points};
  \end{scope}
\end{tikzpicture}\par\vspace{2pt}
{\small\textbf{Figure 2.} The sphere ``unwrapped'' into its grid of angles (drawn with only 4\,$\times$\,8 cells).
Every cell becomes two triangles. In the top and bottom rows one triangle has zero area and is skipped.}
\end{minipage}

\vspace{6pt}\hrule\vspace{2pt}

\begin{multicols}{2}
\step{1. Every point comes from two angles}
Instead of listing points by hand, we describe each point by \emph{where it is}: how far down from the top
({\color{thetac}$\theta$}, from $0$ at the north pole to $\pi$ at the south pole) and how far around
({\color{phic}$\varphi$}, from $0$ to $2\pi$). With $y$ pointing up:
\[
x = r\sin\theta\,\sin\varphi,\qquad y = r\cos\theta,\qquad z = r\sin\theta\,\cos\varphi
\]
Read it off Figure 1: $P$ sits at height $r\cos\theta$, and at distance $r\sin\theta$ from the $y$ axis.
That distance is then split between $x$ and $z$ by the angle $\varphi$.

\step{2. Walk a grid of angles}
Split $\theta$ into \code{LAT\_BANDS} equal steps and $\varphi$ into \code{LON\_BANDS} equal steps.
Each pair $(i, j)$ makes one vertex. The vertices are stored row by row, so vertex $(i,j)$ is at
\[
\text{index}(i,j) = i\cdot(\text{LON\_BANDS}+1) + j .
\]
The last column ($\varphi = 2\pi$) lands on top of the first ($\varphi = 0$). We store it twice anyway,
because that keeps the index math simple.
With $24 \times 48$ bands that is $25 \times 49 = 1225$ vertices.

\step{3. Two triangles per cell}
The four corners of a cell are $a=(i,j)$, $b=(i,j{+}1)$, $c=(i{+}1,j)$, and $d=(i{+}1,j{+}1)$.
The cell (orange in both figures) is cut along the diagonal $a$--$d$ into triangles
\[
(a,\,c,\,d)\quad\text{and}\quad(a,\,d,\,b).
\]
Seen from \emph{outside} the sphere, both triangles go \textbf{counter-clockwise}.
That makes them front faces, so \code{cullMode:\,'back'} throws away only the triangles on the far side.

\step{4. Take care of the poles}
Every vertex in row $i = 0$ is the north pole, so there $a$ and $b$ are the same point.
The triangle $(a,d,b)$ would be a flat line with zero area, so we skip it.
The same thing happens at the south pole ($c = d$), where we skip $(a,c,d)$.
Total: $2\cdot\text{LON\_BANDS}\cdot(\text{LAT\_BANDS}-1) = 2208$ triangles.

\step{5. Normals (and colors) for free}
For a sphere centered at the origin, the outward normal is just the direction from the center to the point:
\[
\mathbf n = \frac{P}{r} = (\sin\theta\sin\varphi,\ \cos\theta,\ \sin\theta\cos\varphi).
\]
We don't need cross products. Neighboring triangles \emph{share} vertices and normals, so the
fragment shader blends the normals smoothly across each triangle. That is why the sphere looks round
even though it is made of flat triangles. (The tetrahedron is different: each face has its own copies of its
vertices, with one flat normal.)
The color is the normal moved from $[-1,1]$ into $[0,1]$: \ $\text{color} = 0.5\,\mathbf n + 0.5$.

\vspace{3pt}
{\small Each vertex is 9 floats: \code{X Y Z | NX NY NZ | R G B}.}
\end{multicols}
\end{document}
